% ------------------------------------------------------------------------------
% 4.3. REQUISITOS ESPECÍFICOS
% ------------------------------------------------------------------------------
\subsection{4.3. Requisitos Específicos}\label{sec:ieee_requisitos_especificos}

\subsubsection{4.3.1. Requisitos Funcionales}\label{sec:ieee_rf}
% [INSTRUCCIÓN]: Especificación formal de cada requisito funcional (RF-01, RF-02, etc.).

A continuación se detallan los requisitos funcionales identificados para el sistema, formulados bajo el esquema de historia de usuario (¿quién?, ¿qué? y ¿por qué?) y acompañados de sus respectivos criterios de aceptación. Las Tablas~\ref{tab:rf01} a~\ref{tab:rf21} presentan cada uno de estos requisitos.

\begin{table}[H]
\caption{Requerimiento funcional RF-01: Registro y autenticación de usuarios}
\label{tab:rf01}
\begin{threeparttable}
\begin{tabular}{>{\raggedright\arraybackslash}p{3.2cm} >{\raggedright\arraybackslash}p{9.3cm}}
\toprule
¿Quién? & Postulante, Administrador DUBSS y Directivo DUBSS. \\
¿Qué? & Quiere registrarse y autenticarse en el sistema desde el portal web o la aplicación móvil. \\
¿Por qué? & Para acceder de forma segura a las funcionalidades correspondientes a su rol, validado mediante el servicio de autenticación de Firebase. \\
\midrule
\multicolumn{2}{l}{\textbf{Criterios de aceptación}} \\
\multicolumn{2}{p{12.5cm}}{
\begin{itemize}
    \item El sistema permite crear una cuenta y autenticarse desde la web y desde la aplicación móvil.
    \item El sistema identifica el rol del usuario autenticado (Postulante, Administrador o Directivo) y lo redirige al módulo correspondiente.
    \item El sistema rechaza credenciales inválidas y muestra un mensaje de error correspondiente.
\end{itemize}
} \\
\bottomrule
\end{tabular}
\begin{tablenotes}
\item \textit{Nota.} Elaboración propia.
\end{tablenotes}
\end{threeparttable}
\end{table}

\begin{table}[H]
\caption{Requerimiento funcional RF-02: Envío de postulación a la Beca Comedor}
\label{tab:rf02}
\begin{threeparttable}
\begin{tabular}{>{\raggedright\arraybackslash}p{3.2cm} >{\raggedright\arraybackslash}p{9.3cm}}
\toprule
¿Quién? & Postulante. \\
¿Qué? & Quiere completar un formulario socioeconómico digital y adjuntar la documentación de respaldo digitalizada, desde el portal web o la aplicación móvil. \\
¿Por qué? & Para iniciar formalmente su proceso de postulación a la Beca Comedor. \\
\midrule
\multicolumn{2}{l}{\textbf{Criterios de aceptación}} \\
\multicolumn{2}{p{12.5cm}}{
\begin{itemize}
    \item El Postulante puede completar la totalidad de los campos del formulario socioeconómico.
    \item El Postulante puede adjuntar uno o más archivos digitales como respaldo documental.
    \item El sistema no permite enviar la postulación si falta algún campo o documento obligatorio.
    \item El sistema confirma la recepción de la postulación y la asocia a la gestión académica vigente.
\end{itemize}
} \\
\bottomrule
\end{tabular}
\begin{tablenotes}
\item \textit{Nota.} Elaboración propia.
\end{tablenotes}
\end{threeparttable}
\end{table}

\begin{table}[H]
\caption{Requerimiento funcional RF-03: Preselección automatizada con Inteligencia Artificial}
\label{tab:rf03}
\begin{threeparttable}
\begin{tabular}{>{\raggedright\arraybackslash}p{3.2cm} >{\raggedright\arraybackslash}p{9.3cm}}
\toprule
¿Quién? & Sistema (módulo de preselección con Inteligencia Artificial). \\
¿Qué? & Quiere precalificar automáticamente cada postulación recibida mediante un modelo de inteligencia artificial que evalúe las variables socioeconómicas y académicas declaradas. \\
¿Por qué? & Para generar un puntaje de vulnerabilidad que permita ordenar objetivamente a los postulantes según su nivel de necesidad, sin intervención manual sobre el total de solicitudes recibidas (hasta 20,000 o más por gestión). \\
\midrule
\multicolumn{2}{l}{\textbf{Criterios de aceptación}} \\
\multicolumn{2}{p{12.5cm}}{
\begin{itemize}
    \item El sistema calcula un puntaje de vulnerabilidad para cada postulación dentro de un tiempo razonable tras su envío.
    \item El sistema procesa automáticamente la totalidad de las postulaciones recibidas, sin requerir revisión manual previa.
    \item El sistema registra el puntaje obtenido y la fecha de cálculo, permitiendo su trazabilidad posterior.
\end{itemize}
} \\
\bottomrule
\end{tabular}
\begin{tablenotes}
\item \textit{Nota.} Elaboración propia.
\end{tablenotes}
\end{threeparttable}
\end{table}

\begin{table}[H]
\caption{Requerimiento funcional RF-04: Panel de monitoreo y toma de decisiones (auditoría y adjudicación)}
\label{tab:rf04}
\begin{threeparttable}
\begin{tabular}{>{\raggedright\arraybackslash}p{3.2cm} >{\raggedright\arraybackslash}p{9.3cm}}
\toprule
¿Quién? & Administrador DUBSS. \\
¿Qué? & Quiere visualizar la lista priorizada de postulantes según su puntaje de vulnerabilidad, auditar los expedientes finalistas y aprobar o rechazar la adjudicación de la beca. \\
¿Por qué? & Para validar los expedientes físicos de los finalistas y generar las actas o comprobantes correspondientes a los beneficiarios aprobados. \\
\midrule
\multicolumn{2}{l}{\textbf{Criterios de aceptación}} \\
\multicolumn{2}{p{12.5cm}}{
\begin{itemize}
    \item El Administrador visualiza la lista de postulantes ordenada de mayor a menor puntaje de vulnerabilidad.
    \item El Administrador puede marcar un expediente como auditado o validado tras revisar el respaldo físico.
    \item El Administrador puede aprobar o rechazar la adjudicación de la beca para cada postulante auditado.
    \item El sistema genera un acta o comprobante digital de adjudicación para cada beneficiario aprobado.
\end{itemize}
} \\
\bottomrule
\end{tabular}
\begin{tablenotes}
\item \textit{Nota.} Elaboración propia.
\end{tablenotes}
\end{threeparttable}
\end{table}

\begin{table}[H]
\caption{Requerimiento funcional RF-05: Registro e impresión de planillas de asistencia}
\label{tab:rf05}
\begin{threeparttable}
\begin{tabular}{>{\raggedright\arraybackslash}p{3.2cm} >{\raggedright\arraybackslash}p{9.3cm}}
\toprule
¿Quién? & Administrador DUBSS y personal del Comedor Universitario. \\
¿Qué? & Quiere generar planillas de asistencia física y digital para su firma y fotocopia de carnet, y registrar digitalmente la asistencia diaria ya verificada de los becados. \\
¿Por qué? & Para mantener un respaldo cruzable frente a las auditorías externas, sin que el sistema dependa de ningún dispositivo de identificación biométrica propio, dado que dicha verificación física es responsabilidad de la empresa concesionaria. \\
\midrule
\multicolumn{2}{l}{\textbf{Criterios de aceptación}} \\
\multicolumn{2}{p{12.5cm}}{
\begin{itemize}
    \item El sistema genera planillas de asistencia imprimibles, listas para firma física y fotocopia de carnet.
    \item El sistema permite registrar digitalmente la asistencia diaria de cada becario, una vez verificada su identidad por el personal del comedor.
    \item El sistema permite cotejar el registro digital de asistencia contra el respaldo físico (planilla firmada) durante una auditoría.
    \item El sistema no depende de ningún dispositivo de identificación biométrica propio; esa verificación es externa al sistema.
\end{itemize}
} \\
\bottomrule
\end{tabular}
\begin{tablenotes}
\item \textit{Nota.} Elaboración propia.
\end{tablenotes}
\end{threeparttable}
\end{table}

\begin{table}[H]
\caption{Requerimiento funcional RF-06: Consultas y notificaciones al postulante}
\label{tab:rf06}
\begin{threeparttable}
\begin{tabular}{>{\raggedright\arraybackslash}p{3.2cm} >{\raggedright\arraybackslash}p{9.3cm}}
\toprule
¿Quién? & Postulante. \\
¿Qué? & Quiere realizar seguimiento en tiempo real del estado de su postulación (recibida, en evaluación, finalista, adjudicada, rechazada) y recibir notificaciones push sobre eventos relevantes. \\
¿Por qué? & Para mantenerse informado del avance y del resultado de su beneficio sin necesidad de acudir presencialmente a la DUBSS. \\
\midrule
\multicolumn{2}{l}{\textbf{Criterios de aceptación}} \\
\multicolumn{2}{p{12.5cm}}{
\begin{itemize}
    \item El Postulante puede consultar en todo momento el estado actual de su postulación.
    \item El sistema notifica automáticamente al Postulante cuando su estado cambia (por ejemplo, al ser seleccionado como finalista o al adjudicársele la beca).
\end{itemize}
} \\
\bottomrule
\end{tabular}
\begin{tablenotes}
\item \textit{Nota.} Elaboración propia.
\end{tablenotes}
\end{threeparttable}
\end{table}

\begin{table}[H]
\caption{Requerimiento funcional RF-07: Dashboard de indicadores de Business Intelligence}
\label{tab:rf07}
\begin{threeparttable}
\begin{tabular}{>{\raggedright\arraybackslash}p{3.2cm} >{\raggedright\arraybackslash}p{9.3cm}}
\toprule
¿Quién? & Directivo DUBSS. \\
¿Qué? & Quiere visualizar un dashboard interactivo de Business Intelligence con indicadores clave del comportamiento histórico de las postulaciones y del uso del beneficio. \\
¿Por qué? & Para apoyar la planificación institucional y la rendición de cuentas con información consolidada de cada gestión. \\
\midrule
\multicolumn{2}{l}{\textbf{Criterios de aceptación}} \\
\multicolumn{2}{p{12.5cm}}{
\begin{itemize}
    \item El Directivo puede visualizar indicadores clave (postulaciones recibidas, becas adjudicadas, asistencia) en un dashboard interactivo.
    \item El sistema permite filtrar los indicadores por gestión académica.
\end{itemize}
} \\
\bottomrule
\end{tabular}
\begin{tablenotes}
\item \textit{Nota.} Elaboración propia.
\end{tablenotes}
\end{threeparttable}
\end{table}

\begin{table}[H]
\caption{Requerimiento funcional RF-08: Gestión de convocatorias y gestiones académicas}
\label{tab:rf08}
\begin{threeparttable}
\begin{tabular}{>{\raggedright\arraybackslash}p{3.2cm} >{\raggedright\arraybackslash}p{9.3cm}}
\toprule
¿Quién? & Administrador DUBSS. \\
¿Qué? & Quiere crear, editar, abrir y cerrar convocatorias de la Beca Comedor, definiendo la gestión académica, las fechas, el número de cupos y los documentos requeridos. \\
¿Por qué? & Para controlar el periodo en que se reciben postulaciones y asociar cada postulación a la gestión vigente. \\
\midrule
\multicolumn{2}{l}{\textbf{Criterios de aceptación}} \\
\multicolumn{2}{p{12.5cm}}{
\begin{itemize}
    \item El Administrador puede registrar una convocatoria con gestión académica, fecha de apertura, fecha de cierre y cupos (por definir con la DUBSS).
    \item El Administrador puede definir la lista de documentos obligatorios de cada convocatoria.
    \item El sistema solo permite enviar postulaciones mientras la convocatoria esté abierta.
    \item Solo puede existir una convocatoria vigente por gestión académica.
\end{itemize}
} \\
\bottomrule
\end{tabular}
\begin{tablenotes}
\item \textit{Nota.} Elaboración propia.
\end{tablenotes}
\end{threeparttable}
\end{table}

\begin{table}[H]
\caption{Requerimiento funcional RF-09: Gestión de usuarios y roles}
\label{tab:rf09}
\begin{threeparttable}
\begin{tabular}{>{\raggedright\arraybackslash}p{3.2cm} >{\raggedright\arraybackslash}p{9.3cm}}
\toprule
¿Quién? & Administrador DUBSS. \\
¿Qué? & Quiere crear, editar, desactivar y asignar roles (Administrador, Directivo y Personal del Comedor) a las cuentas del personal. \\
¿Por qué? & Para garantizar que cada persona acceda únicamente a las funciones de su rol, conforme al control de acceso exigido por el sistema. \\
\midrule
\multicolumn{2}{l}{\textbf{Criterios de aceptación}} \\
\multicolumn{2}{p{12.5cm}}{
\begin{itemize}
    \item El Administrador puede crear y desactivar cuentas de personal y asignarles un rol.
    \item El sistema impide que un usuario acceda a funciones de un rol distinto al asignado.
    \item El sistema conserva el registro de los cambios de rol realizados.
    \item El sistema rechaza el alta de una cuenta con CI o correo electrónico ya registrados.
\end{itemize}
} \\
\bottomrule
\end{tabular}
\begin{tablenotes}
\item \textit{Nota.} Elaboración propia.
\end{tablenotes}
\end{threeparttable}
\end{table}

\begin{table}[H]
\caption{Requerimiento funcional RF-10: Recuperación de contraseña y gestión de perfil}
\label{tab:rf10}
\begin{threeparttable}
\begin{tabular}{>{\raggedright\arraybackslash}p{3.2cm} >{\raggedright\arraybackslash}p{9.3cm}}
\toprule
¿Quién? & Postulante, Administrador DUBSS y Directivo DUBSS. \\
¿Qué? & Quiere recuperar su contraseña en caso de olvido y actualizar sus datos de contacto desde su perfil. \\
¿Por qué? & Para recuperar el acceso a su cuenta sin intervención del personal y mantener actualizados los datos con que se le notifica. \\
\midrule
\multicolumn{2}{l}{\textbf{Criterios de aceptación}} \\
\multicolumn{2}{p{12.5cm}}{
\begin{itemize}
    \item El usuario puede solicitar el restablecimiento de contraseña mediante su correo registrado, a través del servicio de autenticación de Firebase.
    \item El usuario puede modificar sus datos de contacto y su contraseña desde el perfil.
    \item El sistema no permite modificar el CI desde el perfil.
\end{itemize}
} \\
\bottomrule
\end{tabular}
\begin{tablenotes}
\item \textit{Nota.} Elaboración propia.
\end{tablenotes}
\end{threeparttable}
\end{table}

\begin{table}[H]
\caption{Requerimiento funcional RF-11: Digitalización y carga de documentos de respaldo}
\label{tab:rf11}
\begin{threeparttable}
\begin{tabular}{>{\raggedright\arraybackslash}p{3.2cm} >{\raggedright\arraybackslash}p{9.3cm}}
\toprule
¿Quién? & Postulante. \\
¿Qué? & Quiere digitalizar sus documentos con la cámara del teléfono, o cargarlos desde un archivo en la web, y asociarlos a su postulación. \\
¿Por qué? & Para respaldar los datos declarados en el formulario socioeconómico sin presentar carpetas físicas en la etapa inicial. \\
\midrule
\multicolumn{2}{l}{\textbf{Criterios de aceptación}} \\
\multicolumn{2}{p{12.5cm}}{
\begin{itemize}
    \item El Postulante puede capturar un documento con la cámara del dispositivo móvil o adjuntar un archivo desde la web.
    \item El sistema valida el formato y el tamaño del archivo (formatos y tamaño máximo por definir con la DUBSS) y rechaza los no permitidos con un mensaje de error.
    \item El sistema muestra una vista previa y permite reemplazar o eliminar un documento antes de enviar la postulación.
    \item El sistema almacena los archivos en Firebase Storage asociados a la postulación.
\end{itemize}
} \\
\bottomrule
\end{tabular}
\begin{tablenotes}
\item \textit{Nota.} Elaboración propia.
\end{tablenotes}
\end{threeparttable}
\end{table}

\begin{table}[H]
\caption{Requerimiento funcional RF-12: Guardado de borrador de postulación}
\label{tab:rf12}
\begin{threeparttable}
\begin{tabular}{>{\raggedright\arraybackslash}p{3.2cm} >{\raggedright\arraybackslash}p{9.3cm}}
\toprule
¿Quién? & Postulante. \\
¿Qué? & Quiere guardar su postulación incompleta como borrador y retomarla después. \\
¿Por qué? & Para completar el formulario socioeconómico y reunir los documentos sin perder lo ya ingresado. \\
\midrule
\multicolumn{2}{l}{\textbf{Criterios de aceptación}} \\
\multicolumn{2}{p{12.5cm}}{
\begin{itemize}
    \item El Postulante puede guardar el formulario con campos obligatorios sin completar.
    \item El Postulante puede retomar y editar el borrador mientras la convocatoria esté abierta.
    \item Un borrador no se considera enviado ni es procesado por el módulo de preselección.
    \item El sistema descarta o bloquea los borradores no enviados al cerrar la convocatoria.
\end{itemize}
} \\
\bottomrule
\end{tabular}
\begin{tablenotes}
\item \textit{Nota.} Elaboración propia.
\end{tablenotes}
\end{threeparttable}
\end{table}

\begin{table}[H]
\caption{Requerimiento funcional RF-13: Validación de elegibilidad y unicidad de postulaciones}
\label{tab:rf13}
\begin{threeparttable}
\begin{tabular}{>{\raggedright\arraybackslash}p{3.2cm} >{\raggedright\arraybackslash}p{9.3cm}}
\toprule
¿Quién? & Sistema. \\
¿Qué? & Quiere verificar, al recibir una postulación, que el postulante cumple las condiciones de elegibilidad y que no existe otra postulación suya en la misma gestión. \\
¿Por qué? & Para evitar registros duplicados o inválidos que distorsionen la priorización y la adjudicación. \\
\midrule
\multicolumn{2}{l}{\textbf{Criterios de aceptación}} \\
\multicolumn{2}{p{12.5cm}}{
\begin{itemize}
    \item El sistema impide enviar más de una postulación vigente por postulante y por convocatoria.
    \item El sistema verifica las condiciones de elegibilidad (por definir con la DUBSS) antes de aceptar el envío.
    \item El sistema informa al Postulante el motivo cuando su postulación es rechazada por esta validación.
    \item Una postulación devuelta (RF-14) y reenviada no se considera duplicada.
\end{itemize}
} \\
\bottomrule
\end{tabular}
\begin{tablenotes}
\item \textit{Nota.} Elaboración propia.
\end{tablenotes}
\end{threeparttable}
\end{table}

\begin{table}[H]
\caption{Requerimiento funcional RF-14: Devolución de postulación para corrección}
\label{tab:rf14}
\begin{threeparttable}
\begin{tabular}{>{\raggedright\arraybackslash}p{3.2cm} >{\raggedright\arraybackslash}p{9.3cm}}
\toprule
¿Quién? & Administrador DUBSS (y Postulante, como receptor). \\
¿Qué? & Quiere devolver una postulación con observaciones para que el Postulante corrija o complete su documentación y la reenvíe. \\
¿Por qué? & Para evitar descartar postulaciones por defectos subsanables y mantener la trazabilidad de las correcciones. \\
\midrule
\multicolumn{2}{l}{\textbf{Criterios de aceptación}} \\
\multicolumn{2}{p{12.5cm}}{
\begin{itemize}
    \item El Administrador puede devolver una postulación registrando las observaciones.
    \item La postulación cambia al estado devuelta y el Postulante es notificado.
    \item El Postulante puede corregir y reenviar dentro de un plazo (por definir con la DUBSS).
    \item El sistema conserva el historial de devoluciones y reenvíos de la postulación.
\end{itemize}
} \\
\bottomrule
\end{tabular}
\begin{tablenotes}
\item \textit{Nota.} Elaboración propia.
\end{tablenotes}
\end{threeparttable}
\end{table}

\begin{table}[H]
\caption{Requerimiento funcional RF-15: Archivo o rechazo de postulación}
\label{tab:rf15}
\begin{threeparttable}
\begin{tabular}{>{\raggedright\arraybackslash}p{3.2cm} >{\raggedright\arraybackslash}p{9.3cm}}
\toprule
¿Quién? & Administrador DUBSS. \\
¿Qué? & Quiere archivar o rechazar una postulación indicando el motivo. \\
¿Por qué? & Para cerrar formalmente los expedientes que no proceden y dejar constancia de la decisión. \\
\midrule
\multicolumn{2}{l}{\textbf{Criterios de aceptación}} \\
\multicolumn{2}{p{12.5cm}}{
\begin{itemize}
    \item El Administrador debe registrar un motivo para archivar o rechazar.
    \item La postulación pasa al estado correspondiente (archivada o no adjudicada) y deja de aparecer en la lista priorizada activa.
    \item El Postulante puede ver el estado y el motivo registrado.
    \item El sistema conserva la postulación archivada para consulta histórica y reportes.
\end{itemize}
} \\
\bottomrule
\end{tabular}
\begin{tablenotes}
\item \textit{Nota.} Elaboración propia.
\end{tablenotes}
\end{threeparttable}
\end{table}

\begin{table}[H]
\caption{Requerimiento funcional RF-16: Notificación a finalistas, entrega de documentación física y comunicados}
\label{tab:rf16}
\begin{threeparttable}
\begin{tabular}{>{\raggedright\arraybackslash}p{3.2cm} >{\raggedright\arraybackslash}p{9.3cm}}
\toprule
¿Quién? & Administrador DUBSS (y Postulante, como receptor). \\
¿Qué? & Quiere notificar a los finalistas para que presenten su documentación física, registrar la fecha en que se recibe, y publicar comunicados y avisos generales dirigidos a postulantes. \\
¿Por qué? & Para coordinar la auditoría física de los finalistas, dejar constancia de su entrega y mantener informada a la población postulante. \\
\midrule
\multicolumn{2}{l}{\textbf{Criterios de aceptación}} \\
\multicolumn{2}{p{12.5cm}}{
\begin{itemize}
    \item El sistema envía una notificación push a cada finalista indicando el plazo y los documentos que debe presentar.
    \item El Administrador puede registrar la recepción de la documentación física de un finalista con su fecha.
    \item El Administrador puede publicar un comunicado dirigido a todos los postulantes o a un grupo (por ejemplo, los finalistas).
    \item Los comunicados se notifican mediante push y quedan consultables desde la aplicación.
\end{itemize}
} \\
\bottomrule
\end{tabular}
\begin{tablenotes}
\item \textit{Nota.} Elaboración propia.
\end{tablenotes}
\end{threeparttable}
\end{table}

\begin{table}[H]
\caption{Requerimiento funcional RF-17: Publicación de resultados de adjudicación}
\label{tab:rf17}
\begin{threeparttable}
\begin{tabular}{>{\raggedright\arraybackslash}p{3.2cm} >{\raggedright\arraybackslash}p{9.3cm}}
\toprule
¿Quién? & Administrador DUBSS. \\
¿Qué? & Quiere publicar la lista oficial de adjudicados de la gestión. \\
¿Por qué? & Para dar a conocer los resultados de forma transparente y oficial y poner fin al proceso de selección. \\
\midrule
\multicolumn{2}{l}{\textbf{Criterios de aceptación}} \\
\multicolumn{2}{p{12.5cm}}{
\begin{itemize}
    \item El Administrador puede publicar los resultados solo cuando todas las adjudicaciones estén decididas.
    \item La lista publicada incluye a los adjudicados de la gestión y es consultable por los postulantes.
    \item Al publicar, el sistema notifica a cada postulante su resultado.
    \item Una vez publicada, la lista no se modifica sin dejar registro en la bitácora (RF-20).
\end{itemize}
} \\
\bottomrule
\end{tabular}
\begin{tablenotes}
\item \textit{Nota.} Elaboración propia.
\end{tablenotes}
\end{threeparttable}
\end{table}

\begin{table}[H]
\caption{Requerimiento funcional RF-18: Gestión del padrón de becados}
\label{tab:rf18}
\begin{threeparttable}
\begin{tabular}{>{\raggedright\arraybackslash}p{3.2cm} >{\raggedright\arraybackslash}p{9.3cm}}
\toprule
¿Quién? & Administrador DUBSS. \\
¿Qué? & Quiere mantener el padrón de becados vigentes, registrando bajas, suspensiones y reemplazos durante la gestión. \\
¿Por qué? & Para mantener actualizada la lista de beneficiarios activos, de la cual dependen las planillas de asistencia, y reasignar los cupos liberados. \\
\midrule
\multicolumn{2}{l}{\textbf{Criterios de aceptación}} \\
\multicolumn{2}{p{12.5cm}}{
\begin{itemize}
    \item El Administrador puede consultar el padrón de becados vigentes de la gestión.
    \item El Administrador puede dar de baja o suspender a un becado registrando el motivo.
    \item El Administrador puede reemplazar un cupo liberado con el siguiente postulante priorizado.
    \item Los cambios en el padrón se reflejan en las planillas de asistencia generadas posteriormente.
\end{itemize}
} \\
\bottomrule
\end{tabular}
\begin{tablenotes}
\item \textit{Nota.} Elaboración propia.
\end{tablenotes}
\end{threeparttable}
\end{table}

\begin{table}[H]
\caption{Requerimiento funcional RF-19: Historial de asistencia y estado de la beca}
\label{tab:rf19}
\begin{threeparttable}
\begin{tabular}{>{\raggedright\arraybackslash}p{3.2cm} >{\raggedright\arraybackslash}p{9.3cm}}
\toprule
¿Quién? & Becado (Postulante adjudicado). \\
¿Qué? & Quiere consultar su historial de asistencia al comedor y el estado vigente de su beca. \\
¿Por qué? & Para verificar su cumplimiento de los requisitos de asistencia y conservar su beneficio. \\
\midrule
\multicolumn{2}{l}{\textbf{Criterios de aceptación}} \\
\multicolumn{2}{p{12.5cm}}{
\begin{itemize}
    \item El Becado visualiza las fechas de asistencia registradas en la gestión vigente.
    \item El Becado visualiza el estado de su beca (vigente, suspendida o dada de baja) y el motivo, si corresponde.
    \item El Becado solo accede a su propia información.
\end{itemize}
} \\
\bottomrule
\end{tabular}
\begin{tablenotes}
\item \textit{Nota.} Elaboración propia.
\end{tablenotes}
\end{threeparttable}
\end{table}

\begin{table}[H]
\caption{Requerimiento funcional RF-20: Bitácora de acciones}
\label{tab:rf20}
\begin{threeparttable}
\begin{tabular}{>{\raggedright\arraybackslash}p{3.2cm} >{\raggedright\arraybackslash}p{9.3cm}}
\toprule
¿Quién? & Administrador DUBSS. \\
¿Qué? & Quiere consultar un registro de las acciones relevantes realizadas en el sistema (auditoría, devolución, adjudicación, cambios en el padrón y cambios de roles). \\
¿Por qué? & Para garantizar la trazabilidad de las decisiones y facilitar las auditorías sobre el proceso. \\
\midrule
\multicolumn{2}{l}{\textbf{Criterios de aceptación}} \\
\multicolumn{2}{p{12.5cm}}{
\begin{itemize}
    \item El sistema registra automáticamente usuario, acción, expediente afectado, fecha y hora.
    \item El Administrador puede filtrar la bitácora por usuario, tipo de acción y rango de fechas.
    \item Los registros de la bitácora no pueden ser editados ni eliminados desde la interfaz.
\end{itemize}
} \\
\bottomrule
\end{tabular}
\begin{tablenotes}
\item \textit{Nota.} Elaboración propia.
\end{tablenotes}
\end{threeparttable}
\end{table}

\begin{table}[H]
\caption{Requerimiento funcional RF-21: Parametrización del modelo de preselección y reprocesamiento}
\label{tab:rf21}
\begin{threeparttable}
\begin{tabular}{>{\raggedright\arraybackslash}p{3.2cm} >{\raggedright\arraybackslash}p{9.3cm}}
\toprule
¿Quién? & Administrador DUBSS. \\
¿Qué? & Quiere ajustar los criterios y ponderaciones que usa el modelo de preselección y reprocesar los puntajes de las postulaciones de una convocatoria. \\
¿Por qué? & Para calibrar la priorización según los criterios institucionales de vulnerabilidad y corregir puntajes cuando estos cambien. \\
\midrule
\multicolumn{2}{l}{\textbf{Criterios de aceptación}} \\
\multicolumn{2}{p{12.5cm}}{
\begin{itemize}
    \item El Administrador puede consultar y modificar los criterios y ponderaciones del modelo (por definir con la DUBSS).
    \item El Administrador puede solicitar el reprocesamiento de los puntajes de una convocatoria.
    \item El sistema registra la fecha de cada recálculo y conserva el puntaje anterior para trazabilidad (ver RF-03).
    \item El sistema muestra la lista priorizada actualizada tras el reprocesamiento.
\end{itemize}
} \\
\bottomrule
\end{tabular}
\begin{tablenotes}
\item \textit{Nota.} Elaboración propia.
\end{tablenotes}
\end{threeparttable}
\end{table}

\subsubsection{4.3.2. Requisitos No Funcionales}\label{sec:ieee_rnf}
% [INSTRUCCIÓN]: Rendimiento, seguridad, disponibilidad, mantenibilidad y fiabilidad (RNF-01, etc.).

A continuación se detallan los requisitos no funcionales identificados para el sistema, describiendo los atributos de calidad exigidos y sus respectivos criterios de aceptación. Las Tablas~\ref{tab:rnf01} a~\ref{tab:rnf05} presentan cada uno de estos requisitos.

\begin{table}[H]
\caption{Requerimiento no funcional RNF-01: Seguridad y protección de datos personales}
\label{tab:rnf01}
\begin{threeparttable}
\begin{tabular}{>{\raggedright\arraybackslash}p{3.2cm} >{\raggedright\arraybackslash}p{9.3cm}}
\toprule
\multicolumn{2}{l}{\textbf{Descripción}} \\
\multicolumn{2}{p{12.5cm}}{El sistema debe garantizar la confidencialidad, integridad y protección de los datos socioeconómicos y personales de los postulantes, aplicando control de acceso basado en roles y cifrado de la información sensible almacenada y transmitida, en cumplimiento de las políticas institucionales de protección de datos personales de la UAGRM.} \\
\midrule
\multicolumn{2}{l}{\textbf{Criterios de aceptación}} \\
\multicolumn{2}{p{12.5cm}}{
\begin{itemize}
    \item El acceso a los expedientes socioeconómicos está restringido según el rol del usuario autenticado.
    \item La información sensible se transmite y almacena de forma cifrada.
    \item El sistema mantiene un registro (log) de los accesos realizados a cada expediente.
\end{itemize}
} \\
\bottomrule
\end{tabular}
\begin{tablenotes}
\item \textit{Nota.} Elaboración propia.
\end{tablenotes}
\end{threeparttable}
\end{table}

\begin{table}[H]
\caption{Requerimiento no funcional RNF-02: Disponibilidad del servicio}
\label{tab:rnf02}
\begin{threeparttable}
\begin{tabular}{>{\raggedright\arraybackslash}p{3.2cm} >{\raggedright\arraybackslash}p{9.3cm}}
\toprule
\multicolumn{2}{l}{\textbf{Descripción}} \\
\multicolumn{2}{p{12.5cm}}{El sistema debe mantenerse operativo y accesible durante los periodos de alta demanda, particularmente durante las fechas de apertura y cierre de la convocatoria de postulación, minimizando el riesgo de interrupciones del servicio desplegado en AWS.} \\
\midrule
\multicolumn{2}{l}{\textbf{Criterios de aceptación}} \\
\multicolumn{2}{p{12.5cm}}{
\begin{itemize}
    \item El portal web y la aplicación móvil permanecen disponibles durante los picos de postulación.
    \item El sistema cuenta con un mecanismo de recuperación ante caídas del servidor.
\end{itemize}
} \\
\bottomrule
\end{tabular}
\begin{tablenotes}
\item \textit{Nota.} Elaboración propia.
\end{tablenotes}
\end{threeparttable}
\end{table}

\begin{table}[H]
\caption{Requerimiento no funcional RNF-03: Tiempo de respuesta y concurrencia}
\label{tab:rnf03}
\begin{threeparttable}
\begin{tabular}{>{\raggedright\arraybackslash}p{3.2cm} >{\raggedright\arraybackslash}p{9.3cm}}
\toprule
\multicolumn{2}{l}{\textbf{Descripción}} \\
\multicolumn{2}{p{12.5cm}}{El sistema debe procesar la precalificación automática y las operaciones de consulta dentro de tiempos de respuesta aceptables, incluso bajo condiciones de alta concurrencia, considerando el volumen de hasta 20,000 postulaciones por gestión y la revisión masiva simultánea por parte de los auditores.} \\
\midrule
\multicolumn{2}{l}{\textbf{Criterios de aceptación}} \\
\multicolumn{2}{p{12.5cm}}{
\begin{itemize}
    \item El algoritmo de precalificación procesa el total de postulaciones dentro de un tiempo máximo a definir mediante pruebas de carga.
    \item El sistema atiende solicitudes concurrentes de múltiples usuarios sin degradar significativamente su rendimiento.
    \item El sistema maneja correctamente peticiones HTTP duplicadas durante la revisión masiva de carpetas, evitando registros o acciones repetidas (idempotencia).
\end{itemize}
} \\
\bottomrule
\end{tabular}
\begin{tablenotes}
\item \textit{Nota.} Elaboración propia.
\end{tablenotes}
\end{threeparttable}
\end{table}

\begin{table}[H]
\caption{Requerimiento no funcional RNF-04: Usabilidad}
\label{tab:rnf04}
\begin{threeparttable}
\begin{tabular}{>{\raggedright\arraybackslash}p{3.2cm} >{\raggedright\arraybackslash}p{9.3cm}}
\toprule
\multicolumn{2}{l}{\textbf{Descripción}} \\
\multicolumn{2}{p{12.5cm}}{El sistema debe ofrecer interfaces simples e intuitivas, acordes al nivel de conocimiento técnico básico de sus usuarios (Postulante, Administrador y Directivo), sin requerir conocimientos de programación para su uso.} \\
\midrule
\multicolumn{2}{l}{\textbf{Criterios de aceptación}} \\
\multicolumn{2}{p{12.5cm}}{
\begin{itemize}
    \item Un usuario sin experiencia previa en el sistema puede completar el flujo de postulación sin asistencia externa.
    \item El panel de reportes del Directivo se interpreta sin necesidad de conocimientos técnicos.
\end{itemize}
} \\
\bottomrule
\end{tabular}
\begin{tablenotes}
\item \textit{Nota.} Elaboración propia.
\end{tablenotes}
\end{threeparttable}
\end{table}

\begin{table}[H]
\caption{Requerimiento no funcional RNF-05: Compatibilidad y restricciones tecnológicas}
\label{tab:rnf05}
\begin{threeparttable}
\begin{tabular}{>{\raggedright\arraybackslash}p{3.2cm} >{\raggedright\arraybackslash}p{9.3cm}}
\toprule
\multicolumn{2}{l}{\textbf{Descripción}} \\
\multicolumn{2}{p{12.5cm}}{El sistema debe cumplir con las restricciones tecnológicas ya definidas para el proyecto: despliegue en un servidor en la nube de AWS, motor de base de datos PostgreSQL, servicios de autenticación, notificaciones push y almacenamiento de archivos a través de Firebase, y aplicación móvil desarrollada en Flutter con prioridad de compatibilidad para dispositivos Android.} \\
\midrule
\multicolumn{2}{l}{\textbf{Criterios de aceptación}} \\
\multicolumn{2}{p{12.5cm}}{
\begin{itemize}
    \item El backend del sistema se despliega correctamente sobre AWS utilizando PostgreSQL.
    \item La autenticación, las notificaciones push y el almacenamiento de archivos operan a través de Firebase.
    \item La aplicación móvil funciona correctamente en dispositivos Android.
\end{itemize}
} \\
\bottomrule
\end{tabular}
\begin{tablenotes}
\item \textit{Nota.} Elaboración propia.
\end{tablenotes}
\end{threeparttable}
\end{table}

\subsubsection{4.3.3. Requisitos de Interfaz}\label{sec:ieee_interfaces}
% [INSTRUCCIÓN]: Directrices de diseño responsivo, accesibilidad y prototipos de pantalla principales.

Los requisitos de interfaz definen cómo el sistema interactúa con sus usuarios, con el hardware disponible y con otros componentes de software, y se organizan en cuatro categorías.

\begin{itemize}
    \item \textbf{Interfaz de usuario:} el portal web debe desarrollarse con React y TypeScript, y la aplicación móvil con Flutter, ambos bajo un diseño responsivo que se adapte a distintos tamaños de pantalla. La interfaz gráfica de la aplicación móvil debe seguir las directrices de \textit{Material Design}, propias del framework Flutter, garantizando coherencia visual y accesibilidad básica (contraste adecuado, tamaños de fuente legibles) para usuarios sin experiencia técnica previa.
    \item \textbf{Interfaz de hardware:} la aplicación móvil debe ser compatible con dispositivos Android de gama media, como los equipos utilizados para las pruebas del proyecto (Redmi Note 14 Pro y Honor X8), y debe permitir el uso de la cámara del dispositivo para la digitalización de documentos de respaldo.
    \item \textbf{Interfaz de software:} el sistema debe exponer e integrar servicios mediante una API REST desarrollada en Spring Boot, que comunique al frontend web y a la aplicación móvil con la base de datos PostgreSQL y con el módulo de preselección desarrollado en Python. Asimismo, debe integrarse con los servicios de Firebase para la autenticación, el almacenamiento de archivos y el envío de notificaciones push.
    \item \textbf{Interfaz de comunicación:} toda transmisión de datos entre los clientes (web y móvil) y el servidor debe realizarse mediante el protocolo HTTPS, y la autenticación de las peticiones debe validarse mediante JSON Web Tokens (JWT), a fin de proteger la información socioeconómica sensible de los postulantes durante su transmisión.
\end{itemize}

\subsubsection{4.3.4. Requisitos de Base de Datos}\label{sec:ieee_requisitos_bd}
% [INSTRUCCIÓN]: Entidades principales, integridad referencial, respaldos y políticas de retención de datos.

Los requisitos de base de datos definen la estructura y las reglas de gestión de los datos del sistema, sin entrar en el diseño detallado de las tablas, el cual corresponde a una etapa posterior del proyecto.

\begin{itemize}
    \item El sistema debe almacenar los registros de usuarios, postulaciones, expedientes, calificaciones de vulnerabilidad y asistencia en una base de datos relacional PostgreSQL.
    \item El campo correspondiente al carnet de identidad (CI) y al correo electrónico debe ser único dentro de la tabla de usuarios, impidiendo el registro de cuentas duplicadas.
    \item Las tablas de postulaciones, expedientes y asistencia deben mantener integridad referencial mediante claves foráneas hacia la tabla de usuarios, evitando registros huérfanos.
    \item La base de datos debe tener capacidad para almacenar, como mínimo, los registros correspondientes a las 20,000 postulaciones recibidas por gestión académica, conservando el histórico de gestiones anteriores para su explotación en los reportes de Business Intelligence.
\end{itemize}

\subsubsection{4.3.5. Restricciones de Diseño}\label{sec:ieee_restricciones_diseno}
% [INSTRUCCIÓN]: Estándares, lenguajes, herramientas o patrones de arquitectura obligatorios para el diseño del sistema.

El diseño e implementación del sistema está sujeto a las siguientes restricciones técnicas y normativas:

\begin{itemize}
    \item El backend debe desarrollarse en el lenguaje Java utilizando el framework Spring Boot, bajo una arquitectura de servicios RESTful.
    \item El frontend web debe desarrollarse con React y TypeScript, y la aplicación móvil con Flutter y el lenguaje Dart.
    \item El módulo de preselección con Inteligencia Artificial debe desarrollarse en Python, como un componente independiente del backend principal, consumido a través de una API.
    \item El control de versiones del código fuente debe llevarse a cabo en GitHub, y los entornos de desarrollo deben estandarizarse mediante contenedores Docker.
    \item El diseño de la Especificación de Requisitos del Software debe ajustarse al estándar IEEE 830, y el manejo de los datos socioeconómicos de los postulantes debe cumplir con las políticas institucionales de protección de datos personales de la UAGRM.
\end{itemize}

\subsubsection{4.3.6. Requisitos de Rendimiento}\label{sec:ieee_requisitos_rendimiento}
% [INSTRUCCIÓN]: Tiempos de respuesta, capacidad de procesamiento, concurrencia y escalabilidad esperados.

\paragraph{4.3.6.1. Tiempo de respuesta.}

\begin{itemize}
    \item La interfaz principal del portal web y de la aplicación móvil no debe tardar más de 3 segundos en cargar bajo condiciones normales de operación.
    \item El sistema debe mostrar la lista priorizada de postulantes al Administrador DUBSS en un tiempo no mayor a 3 segundos.
    \item El cálculo del puntaje de vulnerabilidad de una postulación individual no debe exceder los 5 segundos desde su envío.
\end{itemize}

\paragraph{4.3.6.2. Capacidad y carga máxima.}

\begin{itemize}
    \item El sistema debe procesar la totalidad de las postulaciones recibidas por gestión (hasta 20,000) sin degradar su rendimiento.
    \item El sistema debe soportar al menos 50 usuarios administrativos concurrentes durante la fase de auditoría de expedientes, sin afectar el tiempo de respuesta.
\end{itemize}

\paragraph{4.3.6.3. Uso de recursos (CPU, memoria y ancho de banda).}

\begin{itemize}
    \item El servidor desplegado en AWS no debe superar el 80\% de uso de CPU ni de memoria durante los picos de carga esperados (apertura y cierre de la convocatoria).
    \item El consumo de ancho de banda de un usuario concurrente no debe degradar de forma perceptible la experiencia de otros usuarios conectados simultáneamente.
\end{itemize}

\paragraph{4.3.6.4. Disponibilidad y tiempo de actividad (\textit{uptime}).}

\begin{itemize}
    \item El sistema debe garantizar una disponibilidad mínima del 99\% durante los periodos de postulación y evaluación de cada gestión académica.
    \item Los mantenimientos o actualizaciones programadas deben realizarse fuera del horario de atención administrativa de la DUBSS.
\end{itemize}

\paragraph{4.3.6.5. Escalabilidad.}

\begin{itemize}
    \item La arquitectura del sistema debe permitir el incremento del volumen de postulaciones en gestiones futuras sin requerir un rediseño estructural.
    \item El backend debe poder escalar mediante la incorporación de instancias adicionales en AWS conforme crezca la demanda del servicio.
\end{itemize}

\paragraph{4.3.6.6. Procesamiento en lote o en tiempo real.}

\begin{itemize}
    \item La autenticación, el envío de postulaciones y la consulta del estado de la beca deben procesarse en tiempo real.
    \item La generación de los reportes y dashboards de Business Intelligence puede procesarse en lotes programados (por ejemplo, de forma diaria o al cierre de la gestión), sin requerir actualización instantánea.
\end{itemize}
